\section{Reducing connected higher-loop subintegrals exactly}
\label{sec:notes-loop-tree-method}
The previous triangle formula reduces independent products through thirteen loops. Connected integrals require identities that account for their shared loop momenta. We therefore integrate a massless two-line subgraph whenever it belongs to a supported connected tree attached to a one-scale triangle. Analytic integration of massless insertions is an established part of parametric reduction~\cite{Ruijl2017Forcer}. Appendix~\ref{app:notes-loop-tree} derives the identity for the connected vertex family used here, including its signs, scale dependence, and allowed powers.

Each bubble integration removes one loop and transfers a generally noninteger power to its parent line. Repeating this step from the leaves towards the root leaves a triangle whose first power depends on $d$. A Feynman-parameter integral evaluates this triangle. Dividing by the same connected integral with unit physical powers cancels the common Gamma functions, and finite Gamma shifts give a rational coefficient in $d$ and the nonzero external scale. The implementation keeps this coefficient in factored form, which avoids expanding polynomials that would later cancel.

The new family contains stars, chains, and branched trees with one through thirteen loops. All internal lines are massless, two external momenta are lightlike, and their sum has nonzero squared momentum. Positive powers can vary on every physical line. The middle root-triangle power can also be zero or negative, while the auxiliary scalar-product powers remain zero. These conditions define a connected family with an explicit reduction proof. A generic massive insertion or multi-scale root graph needs further identities.

We apply the identity before loading the remaining Kira reduction workflow. If the denominators, kinematics, or powers fall outside its domain, the application uses the preceding exact methods and trained native choices. We retain the existing batch and sparse-seed models without fitting them to the new test results. The added speedup therefore comes from exact integration and its early application. A separate fresh two-loop comparison measures whether the retained learned choices still reduce the remaining native time.
